\subsection{Prediction-Based Scheduling}
Prediction-based scheduling tries to estimate output length ahead of time, to cut down on head-of-line blocking. Fu et al. \cite{fu2024ltr} do this a bit differently: instead of predicting the exact number of output tokens, they just rank requests by relative output length. They train an OPT-based predictor with ListMLE. To check how good the ranking is, they compare it against Shortest-Job-First (SJF) and Shortest-Remaining-Time-First (SRTF) using Kendall's Tau. Once integrated with vLLM, P90 latency drops by up to 2.8$\times$ for chatbot serving, and throughput goes up by as much as 6.5$\times$ for synthetic data generation. The approach also promotes requests that already waited past a threshold, which helps keep starvation down.

The predictor does not hold up the same way across every dataset, though. A model trained on ShareGPT gets a Kendall's Tau of 0.54 on ShareGPT, but only 0.45 on LMSYS-Chat-1M. The same drop shows up going the other way: a model trained on LMSYS-Chat-1M falls from 0.62 to 0.40 when tested on ShareGPT.

Saravana Kumar et al. \cite{kumar2024latency} take a different combination. They pair an OPT-125M ranking predictor with vLLM's swap-based preemption, and under high request load, this cuts latency by up to 2.1$\times$ compared to FCFS. Their Fig. 3 shows something worth noting too: the improvement gets weaker on a prompt subset that was not part of training. That is part of what pushed us to run our own out-of-distribution (OOD) evaluation. Park and Kwon \cite{park2024prefixoutput} take yet another route, building on PREBLE's per-GPU load estimate but adding output length into it. On an eight-GPU system, this gets latency reductions of 14.31\% at 64 requests/s and 28.89\% at 128 requests/s. Their main experiments do lean on a perfect oracle, though, one that always knows the true output length in advance.



\subsection{Iteration-Level and Memory-Aware Serving Systems}
ORCA \cite{yu2022orca} came up with iteration-level scheduling, paired with selective batching. After every decoding step, it rebuilds the active batch, so finished requests can leave right away and new ones can join without waiting for the whole batch to wrap up. On GPT-3 175B, ORCA gets 36.9$\times$ more throughput than FasterTransformer, and latency stays about the same. Still, ORCA sticks to an iteration-level FCFS order based on when requests arrive. It does not look at predicted request size or how uncertain that prediction might be.

Sarathi-Serve \cite{agrawal2024sarathiserve} pushes the throughput-latency tradeoff further with chunked prefills and stall-free batching. This cuts down on stalls when long prefills end up sharing a batch with decoding that is already running. But its focus is batching efficiency, not starvation caused by uncertain output lengths. Llumnix \cite{sun2024llumnix} works on load balance and tail latency instead, by moving live requests across model instances. Our simulator only looks at one model instance, so this is outside its scope. LoongServe \cite{jiao2024loongserve} takes yet another angle: elastic sequence parallelism, which adjusts GPU allocation across requests and across the different phases of inference. For long-context workloads, it gets up to 3.85$\times$ more throughput than chunked-prefill and prefill-decode disaggregation baselines.

Most of these systems are really about batching, moving requests around, or running things in parallel. What we look at is different and narrower: how to order a single request queue when nobody knows the real request sizes ahead of time.



\subsection{Classical Scheduling Theory and Simulation Frameworks}

The aging term in our priority formula comes directly from an old idea in scheduling theory. Brinch Hansen proposed response-ratio scheduling, where request priority grows the longer a task waits in queue~\cite{hansen1973os}. This prevents a request from staying behind new arrivals forever, because wait time continuously increases its priority score. The beta term was built with this same logic in mind, adapted for systems where only a predicted request size is available instead of the real one. SARATHI also talks about this same tension between speed and fairness, confirming that adding aging mechanisms is a common answer to starvation problems~\cite{agrawal2023sarathi} and not something specific of the design we chose to pursue in our project.

For the simulation setup, existing tools were analyzed before making the choice of building a new one. Vidur~\cite{agrawal2024vidur}, LLMServingSim~\cite{cho2024llmservingsim}, and Frontier~\cite{feng2026frontier} all share the same core idea of a discrete-event structure where requests arrive over time and a scheduler selects the next task, which consolidated our idea to use the same approach. However, those simulators focus heavily on hardware details, profiling real GPU kernels to estimate execution times with high precision. That level of detail was not something this project needed, since the goal here was comparing scheduling policies, not modeling GPU performance itself. What actually mattered most for the decision to build our own tool differently than the ones that we found in the literature was having full control over one specific metric, separating requests that complete from requests that get cancelled after waiting too long, and keeping the two apart in every calculation we run.